\documentclass{standalone}

\usepackage[margin=1.0in]{geometry}
\usepackage{tikz}
\usetikzlibrary{arrows.meta, shapes.geometric, positioning}

\begin{document}
% Custom TikZ styles with high-contrast colors
\tikzset{
  process/.style={
    trapezium,
    trapezium left angle=75,
    trapezium right angle=105,
    draw=teal!80!black,
    fill=cyan!30!white,
    minimum height=3.5em,
    text width=8em,
    font=\sffamily\bfseries\color{black},
    align=center,
    node distance=2.2cm
  },
  evaluate/.style={
    ellipse,
    draw=blue!80!black,
    fill=blue!10!white,
    minimum height=2.5em,
    text width=7em,
    font=\sffamily\bfseries\color{black},
    align=center,
    node distance=2.2cm
  },
  terminal/.style={
    rectangle,
    draw=red!80!black,
    fill=red!10!white,
    rounded corners=2pt,
    minimum height=3em,
    text width=6em,
    font=\sffamily\bfseries\color{black},
    align=center,
    node distance=3cm
  },
  link/.style={
    very thick,
    color=teal!90!black,
    -{Stealth[scale=1.2]}
  },
  feedback/.style={
    thick,
    dashed,
    color=blue!90!black,
    -{Stealth[scale=1.1]}
  }
}

\begin{tikzpicture}[auto]
  % Nodes (renamed for Computer Science domain)
  \node [process] (acquisition_node) {Data Acquisition};
  \node [evaluate, below=of acquisition_node] (preprocess_node) {Preprocessing Stage};
  \node [process, below=of preprocess_node] (feature_node) {Feature Extraction};
  \node [process, below=of feature_node] (training_node) {Model Training};
  \node [process, below=of training_node] (validation_node) {Model Validation};
  \node [process, left=3cm of feature_node] (iteration_node) {Iteration Check};
  \node [terminal, below=of validation_node] (terminal_node) {Process Complete};

  % Invisible node for feedback loop
  \node [draw=none, right=2.5cm of validation_node] (feedback_anchor) {};

  % Edges (with new color and arrow styles)
  \draw[link] (acquisition_node) -- (preprocess_node);
  \draw[link] (preprocess_node) -- (feature_node);
  \draw[link] (feature_node) -- (training_node);
  \draw[link] (training_node) -- (validation_node);
  \draw[link] (validation_node) -- (terminal_node);
  \draw[feedback] (validation_node) -| (feedback_anchor.center);
  \draw[feedback] (feedback_anchor.center) |- (feature_node);
  \draw[link] (feature_node) -- (iteration_node);
  \draw[link] (iteration_node) -- (terminal_node);

\end{tikzpicture}
\end{document}
